\subsection{Problem Definition}
\label{sec:method_problem_definition}

Building upon the overall architectural framework introduced in Section~\ref{sec:method_framework}, we formalize the zero-shot fault diagnosis task as a cross-modal generative problem. In a typical industrial monitoring scenario, we consider a dataset comprising labeled vibration signals from a set of seen fault classes, denoted as $\mathcal{S} = \{s_1, s_2, \dots, s_K\}$, where $K$ is the number of observed categories. For these classes, we have a training set $\mathcal{D}_{tr} = \{(x_i, y_i, a_i) \mid x_i \in \mathcal{X}, y_i \in \mathcal{S}, a_i \in \mathcal{A}\}$, where $x_i$ represents the raw vibration feature vector, $y_i$ is the corresponding class label, and $a_i$ is the semantic descriptor.

The core challenge addressed by the CESM-Diff model is the presence of a disjoint set of unseen fault classes, denoted as $\mathcal{U} = \{u_1, u_2, \dots, u_L\}$, where $\mathcal{S} \cap \mathcal{U} = \emptyset$. For these $L$ unseen classes, no physical vibration samples are available during the training phase. However, we assume access to their corresponding semantic descriptors $\mathcal{A}_{\mathcal{U}}$. While traditional zero-shot learning (ZSL) methods often rely on handcrafted binary attribute vectors—which are prone to human bias and lack representational depth—our approach utilizes the high-dimensional continuous latent space of the Contrastive Language-Image Pretraining (CLIP) model. Specifically, each class $c \in \mathcal{S} \cup \mathcal{U}$ is mapped to a semantic embedding vector $\mathbf{z}_c \in \mathbb{R}^d$ via a frozen CLIP text encoder, providing a rich, multi-modal foundation for class-level knowledge transfer.

The primary goal of our framework is to learn a conditional generative mapping $G: \mathcal{A} \to \mathcal{X}$ that can synthesize discriminative vibration features $\hat{\mathbf{x}}$ for the unseen classes $\mathcal{U}$ given their semantic embeddings $\mathbf{z}_u$. By training on the labeled samples from $\mathcal{S}$, the model must capture the underlying topological relationship between the semantic space $\mathcal{Z}$ and the vibration manifold $\mathcal{X}$. Once the generator $G$ is optimized, it is used to produce a synthetic dataset $\hat{\mathcal{D}}_{unseen} = \{(\hat{x}_j, y_j) \mid y_j \in \mathcal{U}\}$ of sufficient quality and diversity to train a robust downstream classifier capable of identifying novel fault modes in real-time sensor streams.
